\documentclass[11pt]{article}
\usepackage[a4paper,margin=1in]{geometry}
\usepackage[T1]{fontenc}
\usepackage{lmodern,microtype}
\usepackage{amsmath,amssymb,amsthm,mathtools}
\usepackage{booktabs,xcolor}
\usepackage[numbers,sort&compress]{natbib}
\usepackage[colorlinks=true,linkcolor=blue!55!black,citecolor=blue!55!black]{hyperref}
\newtheorem{theorem}{Theorem}[section]
\newtheorem{proposition}[theorem]{Proposition}
\newtheorem{corollary}[theorem]{Corollary}
\theoremstyle{remark}\newtheorem{remark}[theorem]{Remark}
\title{A Direct Factorwise Deterministic-Numerator Tail Certificate}
\author{Bin Wang}\date{July 2026}
\begin{document}
\maketitle

\begin{abstract}
We derive a direct all-order target-tail bound from the parity-resolved
factorization of RH-263.  At the unit disk and first omitted order $29$, the
Fredholm trace ideal, the two even endpoint factors, and the exact odd factor
give a certified logarithmic budget below $2.6624745\times10^{-5}$.  The
relative exponential error is below $2.6625100\times10^{-5}$.  The argument
uses no angular sampling and does not assert a cloud coefficient bridge or a
uniform quotient theorem.
\end{abstract}

\section{Tail decomposition}
Let $u=z/r_H$, $r_H=17/20$, $W=|u|^2$, and $t=|u|/\lambda$.
The RH-15 factorization is
\[
 \log G(u)=P_1u+\log\det(I-u^2T)+\log A_*(u^2)+\log B(u^2)+\log C(u).
\]
At $|z|=1$, put $W=(20/17)^2$ and $t=20/(17\lambda)$.  The RH-13
certificate supplies a nuclear bound $\nu_1$ and bounds $p_1\ge\|T\|$,
$p_2\ge\|T^2\|$, $q\ge\|T^3\|$ as in RH-262 \citep{WangRH13}.

For $n=3k+j$, $j\in\{1,2,3\}$, the trace ideal gives
\[
 \nu(T^n)\le q^k\nu_j,
 \qquad \nu_2=p_1\nu_1,\quad \nu_3=p_2\nu_1.
\]
If $n\ge n_0$, then
\begin{equation}
 \sum_{n\ge n_0}\frac{|\operatorname{tr}T^n|W^n}{n}
 \le \sum_{j=1}^3
 \frac{\nu_jW^j(qW^3)^{k_j}}
 {(3k_j+j)(1-qW^3)},
 \label{eq:fredtail}
\end{equation}
where $k_j$ is the first integer with $3k_j+j\ge n_0$.

\begin{proposition}[Endpoint tails]
For $n_0\ge1$ and $y=t^2<1$,
\[
 \sum_{n\ge n_0}\frac{|[u^{2n}]\log A_*|W^n}{1}
 \le\frac{y^{n_0}}{n_0(1-y)},
\]
and the corresponding $B$ bound is this quantity divided by
$2(1-\lambda^{-2})$.  For odd order $m\ge3$ the exact coefficient is
$t^m/(1+\lambda^{-m})$ before the $1/m$ logarithmic normalization.
\end{proposition}

\begin{proof}
Use $|\lambda^{-2n}/(1+\lambda^{-n})|\le\lambda^{-2n}$ for $A_*$,
$d_n\le\lambda^{-2n}/(1-\lambda^{-2})$ for $B$, and sum a geometric
series after replacing each denominator by its first denominator.  The odd
statement is the exact RH-263 identity.
\end{proof}

\section{Certified order-29 budget}
The first omitted even trace power is $n_0=15$ and the first omitted odd order
is $29$.  The 100-, 150-, and 200-decimal Arb replays all satisfy the strict
outward comparisons
\begin{center}
\begin{tabular}{lr}
\toprule
component & certified upper bound\\
\midrule
Fredholm even tail & $0.000019045786$\\
$A_*$ tail & $0.000003066234$\\
$B$ tail & $0.000002376597$\\
even total & $0.000024488616$\\
odd endpoint tail & $0.000002136130$\\
total logarithmic tail & $0.000026624745$\\
relative exponential error & $0.000026625100$\\
\bottomrule
\end{tabular}
\end{center}

\begin{theorem}[Direct factorwise target tail]\label{thm:tail}
For the deterministic Hardy-scaled numerator at $|z|\le1$,
\[
 \sum_{n\ge29}\frac{|a_n|}{n}<0.000026624745,
 \qquad
 \exp\!\left(\sum_{n\ge29}\frac{|a_n|}{n}\right)-1
 <0.000026625100.
\]
\end{theorem}

\begin{proof}
Orders $n\ge29$ split into even $n=2m$ with $m\ge15$ and odd orders.
Apply \eqref{eq:fredtail} to the Fredholm logarithm, the proposition to
$A_*$ and $B$, and the exact odd series to $C$ after the $P_1=b_1$
cancellation.  The displayed outward endpoints sum to the stated bounds.
\end{proof}

The result is deterministic and all-order in the sense of the analytic
factorization; it is not an all-order theorem about noisy selected clouds.
The five-component ledger remains $(0,0,0,1,1)$: no legal anchored head,
cloud coefficient bridge, or uniform quotient tail has been supplied.  Gates
A--E remain false/open, and no Hilbert--Polya or zeta conclusion is made.

\bibliographystyle{plainnat}\bibliography{references}
\end{document}
